\section{Finite canonical histories identify the original global terminal manifold}

The local-to-global comparison in
\cite{bierstone-milman08-functoriality,wlodarczyk05-hironaka}
can be made on the already constructed original terminal manifold whenever
an actual finite canonical terminal history is available. The finite history
is genuine geometric input; no local comparison or compatibility datum is
assumed. For compact original bases the canonical driver supplies this
history internally. The construction below does not assert existence of a
finite history on an arbitrary noncompact base.

\begin{definition}[Actual local comparison with a finite terminal history]
\label{hir:original-global-terminal-finite-local-comparison}
\uses{hir:original-terminal-open-actual-history-data,hir:original-input-covering-local-canonical-terminal-atlas,hir:original-independent-local-terminal-overlap-unique-comparison}
\lean{Hironaka.HironakaOriginalTerminalFiniteIdentification.finiteTerminalLocalComparison,Hironaka.HironakaOriginalTerminalFiniteIdentification.finiteTerminalLocalComparison_proj}
\leanok
\source{bierstone-milman08-functoriality,wlodarczyk05-hironaka}
Let $M$ be an original marked ideal in the invariant domain on $X$, and
let $T:D\to X$ be an actual finite canonical terminal history. For every
original atlas open $U_i$, construct a biholomorphism
$c_i:D_i\to p_T^{-1}U_i$ from the original local final space. It satisfies
$p_T(c_i(y))=p_i(y)$ at every point.
\end{definition}
\begin{proof}
Restrict the actual history $T$ to the entire open $X$. Its base is the
literal whole-open subtype and its final space is the literal whole
projection preimage. Both are canonically biholomorphic to the original
spaces by subtype inclusion. The original restriction theorem supplies
canonicality and terminality. Compare this actual whole-open history with
the original history on $U_i$, over the common open $U_i$. The common-open
comparison theorem supplies a genuine biholomorphism, including both
holomorphic directions and the whole projection equality. The local
history's preimage of its entire base is canonically its whole final space;
the whole-open history's preimage of $U_i$ is canonically the literal
$p_T^{-1}U_i$ by forgetting the redundant whole-open subtype. Composing
these actual subtype identifications gives $c_i$ and the stated equality.
\end{proof}

\begin{lemma}[Finite local comparisons obey the original gluing relation]
\label{hir:original-global-terminal-finite-local-compatibility}
\uses{hir:original-global-terminal-finite-local-comparison,hir:original-independent-local-pairwise-terminal-triple-transition-laws,hir:original-terminal-overlap-further-restriction-intrinsic-naturality}
\lean{Hironaka.HironakaOriginalTerminalFiniteIdentification.finiteTerminalLocalComparison_transition}
\leanok
\source{bierstone-milman08-functoriality,wlodarczyk05-hironaka}
On every actual original overlap, the finite local comparison after the
original transition equals the first finite local comparison, as an equality
of points of $D$.
\end{lemma}
\begin{proof}
Use the original triple-transition theorem on the histories over $U_i$,
$U_j$, and the entire open $X$, with the third history the actual restriction
of $T$ to $X$. Every point of the first overlap also lies over $X$, so the
required third-open membership is automatic. That theorem identifies the
first-to-third comparison with the composite first-to-second-to-third
comparison. The smaller-open comparison theorem identifies these actual
comparisons with the restrictions of the comparisons defining $c_i$ and
$c_j$. Forget the redundant literal whole-open subtypes. The resulting
equality is precisely the claimed identity in $D$.
\end{proof}

\begin{theorem}[The original global terminal manifold is biholomorphic to a finite terminal history]
\label{hir:original-global-terminal-finite-whole-biholomorphism}
\uses{hir:original-global-terminal-finite-local-compatibility,hir:original-terminal-local-actual-saturated-range-biholomorphism,hir:original-input-local-terminal-glued-whole-continuous-projection}
\lean{Hironaka.HironakaOriginalTerminalFiniteIdentification.originalTerminalFiniteBiholomorphism,Hironaka.HironakaOriginalTerminalFiniteIdentification.originalTerminalFiniteBiholomorphism_proj}
\leanok
\source{bierstone-milman08-functoriality,wlodarczyk05-hironaka}
For the same original input and actual finite canonical terminal history,
there is a genuine biholomorphism $C:\widetilde X_M\to D$ from the original
global terminal manifold. It equals $c_i$ on every original local inclusion,
and its entire projection is the same original glued projection:
$p_T C=P_M$.
\end{theorem}
\begin{proof}
The preceding compatibility gives a continuous map from the actual
topological gluing to $D$. Its equality with each $c_i$ proves the whole
projection identity, since the original inclusions are jointly surjective.
For surjectivity choose, for a point $d\in D$, the original atlas index at
$p_T(d)$; the corresponding open contains that base point, so $d$ belongs
to $p_T^{-1}U_i$ and the inverse of $c_i$ supplies a preimage.
For injectivity, equal images have equal base projections. The saturation
identity therefore puts both original global points in the same original
local range. Use that range's actual local inverse representatives. Their
images under the same injective $c_i$ coincide, so the representatives and
then the global points coincide.

To prove forward holomorphicity at a global point, use an original saturated
range containing it. On that genuine open neighborhood the glued map is
$c_i$ composed with the actual local range inverse. Both are holomorphic.
For inverse holomorphicity at $d$, use the true open neighborhood
$p_T^{-1}U_i$ selected above. There the inverse equals the original local
inclusion composed with $c_i^{-1}$. Both are holomorphic. Equality on these
actual open neighborhoods transfers the regularity to the whole maps.
Bundle bijectivity and the two proved holomorphic directions. These
arguments also include empty manifolds and dimension zero.
\end{proof}
